\section{Results}
\label{sec:results}

% (~8-10 pages)

\subsection{Dataset Characteristics and Preprocessing}
% Sample counts after QC for HMP, AGP, and curatedMetagenomicData
% NT pathway annotation: genera mapped per pathway
% NT production score distributions across cohorts
% Table 1: Dataset summary statistics

\subsection{Co-occurrence Network Structure and NT Subnetworks}
\label{subsec:network_results}
% Full network topology across health states
% NT-specific subnetwork comparisons:
%   - Serotonin subnetwork: Enterococcus-Clostridium-Ruminococcus structure
%   - GABA subnetwork: Lactobacillus-Bifidobacterium-Bacteroides connectivity
%   - Dopamine subnetwork: smaller, Enterococcus-Escherichia core
%   - SCFA subnetwork: Faecalibacterium-Roseburia-Eubacterium hub structure
% Figure: NT subnetwork visualizations

\subsection{Persistence Diagrams Reveal Topological Signatures of NT Production}
\label{subsec:persistence_results}
% Full network: H$_0$, H$_1$, H$_2$ features
% NT subnetworks: pathway-specific topological profiles
%   - High-serotonin samples: more persistent H$_1$ loops in serotonin subnetwork
%   - Cross-feeding loops between SCFA and serotonin subnetworks
% Betti curves and persistence landscapes as NT-state discriminators
% Figure: Persistence diagrams, Betti curves for high vs. low NT groups

\subsection{Mycobiome Disruption Fragments NT-Network Topology}
\label{subsec:mycobiome_results}
% High-fungal vs. low-fungal burden comparison:
%   - Loss of H$_1$ features in GABA subnetwork (Candida displaces Lactobacillus)
%   - Fragmentation of H$_0$ in serotonin subnetwork
%   - Disruption scores correlate with topological degradation
% Per-fungal-genus effects: Candida vs. Aspergillus vs. Malassezia
% Figure: Topological comparison high vs. low fungal burden

\subsection{AMR Carriers as Topological Disruptors}
\label{subsec:amr_results}
% AMR burden correlates with NT-network topology metrics
% Resilience simulation: topological degradation curve
%   - NT cross-feeding loops collapse at ~30-40\% non-AMR removal
%   - ``Topological tipping point'' for each NT pathway
% HGT edges: high-HGT AMR carriers form a connected subcomponent
%   that does not participate in NT cross-feeding
% Figure: Resilience curves showing topological degradation

\subsection{Topological Regime Transitions in NT Networks}
\label{subsec:regime_results}
% Sliding window analysis ordered by NT production gradient
% Topological transitions precede compositional shifts
% Figure: Multi-panel regime detection with topological features

\subsection{Topological Features Predict NT Production and Mood Outcomes}
\label{subsec:prediction_results}
% Correlation with NT production scores
% Correlation with AGP self-reported mental health outcomes
% Cross-cohort validation with curatedMetagenomicData
% ML classification: topological features vs. diversity metrics
% Table: Classification performance comparison
